\documentclass[dvipdfmx,a4paper,10pt,twocolumn]{jsarticle}

% 余白
\usepackage[top=13mm,bottom=13mm,left=12mm,right=12mm]{geometry}

% 数式
\usepackage{amsmath,amsfonts}
\usepackage{bm}

% 画像
\usepackage{graphicx}
\graphicspath{{results/note/}}
\usepackage[haranoaji]{pxchfon}

% 表
\usepackage{array}

\setlength{\columnsep}{7mm}
\setlength{\parindent}{1em}
\setlength{\parskip}{0pt}

\begin{document}

\title{\vspace{-2em}
M{\o}lmer--S{\o}rensenゲートの解析
\vspace{-1em}}
\author{平井 希空}
\date{\today}
\maketitle
\vspace{-3em}

\section{Motivation}

平均ゲートinfidelityは性能比較には便利だが，コヒーレントな較正ずれと確率的な散逸を区別できない．特にM{\o}lmer--S{\o}rensen（MS）ゲートでは，運動自由度を消去した後に相関二量子ビット誤差が生じ得る．そこで，master equation（ME）から得た量子チャネルをQPTで再構成し，理想MSゲートに対する\textbf{error channel}を用いて，実験で調整可能な誤差成分へ分解する．

\section{Prior work \& Method}

MSゲートは二色光により運動モードを介して
$U_{\rm MS}=\exp(i\pi XX/4)$を実装する\cite{MS}．本計算ではmotional heating，motional dephasing，spin dephasing，photon scatteringを含むMEを解き，16個の入力演算子からQPTチャネル$\mathcal E$を構成する\cite{QPT}．理想操作を除いた
\begin{equation}
 \mathcal E_{\rm err}=\mathcal U_{\rm MS}^{\dagger}\circ\mathcal E
 \label{eq:error_channel}
\end{equation}
をCPTPチャネルへ射影し，Pauli transfer matrix $R_{\rm err}$に変換する．一ゲート当たりのeffective generatorを
\begin{equation}
 K=\log R_{\rm err}
 \end{equation}
とし，Pauli基底で
\begin{equation}
 K\simeq\sum_{P\ne I}h_P\mathcal H_P
       +\sum_{P\ne I}\gamma_P\mathcal D_P,
 \quad
 \begin{cases}
 \mathcal H_P(\rho)=-i[P,\rho],\\
 \mathcal D_P(\rho)=P\rho P-\rho
 \end{cases}
 \label{eq:generator}
\end{equation}
と分解する\cite{GST}．$h_{XX}$はcoherentなXX回転誤差，$\gamma_{XX}$はstochasticな相関Pauli-XX誤差を表す．散逸係数は非負最小二乗法で推定した．

\newpage
\section{Results}

\begin{center}
\small
\setlength{\tabcolsep}{3.4pt}
\renewcommand{\arraystretch}{1.08}
\begin{tabular}{c|c|c|c|c}
\hline
$\bar n$ & $1-F_{\rm avg}$ & $h_{XX}$ & $\gamma_{XX}$ & $\gamma_{XI,IX}$ \\
 &  & [rad/gate] & [/gate] & [/gate] \\
\hline
$0.01$ & $1.95\!\times\!10^{-3}$ & $1.37\!\times\!10^{-2}$ & $1.64\!\times\!10^{-4}$ & $5.51\!\times\!10^{-4}$ \\
$1$    & $3.64\!\times\!10^{-3}$ & $2.88\!\times\!10^{-2}$ & $6.16\!\times\!10^{-4}$ & $1.07\!\times\!10^{-3}$ \\
$4$    & $1.23\!\times\!10^{-2}$ & $7.27\!\times\!10^{-2}$ & $4.29\!\times\!10^{-3}$ & $2.48\!\times\!10^{-3}$ \\
\hline
\end{tabular}
\end{center}

$\bar n$の増加とともにinfidelity，$h_{XX}$，$\gamma_{XX}$はいずれも増大した．低温では局所的な$XI/IX$散逸が$XX$散逸より大きいが，$\bar n=4$では$\gamma_{XX}$が最大のPauli散逸成分となる．すなわち，熱占有数の増加により，誤差チャネルは\textbf{局所誤差優勢から相関XX誤差優勢へクロスオーバーする}．実験的には，$h_{XX}$はパルス面積・detuningの再較正で補償可能である一方，$\gamma_{XX}$はユニタリ補正では除けず，冷却および運動コヒーレンスの改善が必要となる．

ここで$\gamma_{XX}$は\emph{有効二量子ビットチャネル中の相関誤差}であり，異なる環境浴の統計相関を直接意味しない．浴相関を主張するには，MEに非対角Kossakowski成分または相互相関スペクトルを導入し，独立浴モデルと比較する必要がある．

\begin{thebibliography}{99}
\small

\bibitem{MS}
A. S{\o}rensen and K. M{\o}lmer,
Phys. Rev. A \textbf{62}, 022311 (2000).

\bibitem{QPT}
I. L. Chuang and M. A. Nielsen,
J. Mod. Opt. \textbf{44}, 2455--2467 (1997).

\bibitem{GST}
E. Nielsen \textit{et al.},
Quantum \textbf{5}, 557 (2021).

\end{thebibliography}

\end{document}
